\section{Definição de Variável Aleatória Discreta}

Uma variável aleatória discreta é aquela cuja imagem é um conjunto que pode ser colocado
em correspondência de 1 para 1 com os naturais, conjuntos esses chamados de
\textit{enumeráveis}.

\begin{def:variavel discreta}
Seja $X$ uma variável aleatória. A variável $X$ é discreta caso $\im{X}$ seja um conjunto
enumerável.
\end{def:variavel discreta}

Para as v.a.'s discretas, a função de probabilidade modela a própria probabilidade
do valor no exeperimento, seja com a interpretação da frequência com a qual ele se
repetiu ou da crença de que ele será o resultado do experimento. A função é então chamada
de \textbf{função de massa de probabilidade} (f.m.p.).

\begin{def:fmp}
Seja $X: \Omega \to \mathbb{R}$ uma v.a. discreta. A função de massa de probabilidade
$p_X: \mathbb{R} \to [0, 1]$ associa todo $x \in \mathbb{R}$ a uma probabilidade.
\end{def:fmp}

As realizações possíveis de $X$, isto é, que possuem probabilidade estritamente maior
que zero, formam o suporte de $X$.

\begin{def:suporte discreta}
Seja $X: \Omega \to \mathbb{R}$ uma v.a. discreta. O suporte de $X$, denotado por $\Omega_X$,
é o conjunto de todos os valores em $\mathbb{R}$ com probabilidade não nula.
\[
	\Omega_X = \{x \in \mathbb{R} \mid p_x(x) > 0 \}
\]
\end{def:suporte discreta}

Uma vez que $\Omega_X \subset \im{X}$, então o suporte de um v.a. discreta é também
enumerável. Além disso, é natural pensarmos que, se $\Omega_X$ compreende todos os valores
ossíveis para $X$, então a probabilidade de que $X$ tenha realização no seu suporte deve
er igual a 1.

\begin{prop:fmp suporte}
Se $X$ é uma v.a. discreta com f.m.p. igual a $p_X$, então
\[
	\sum_{x \in \Omega_X} p_X(x) = 1
\]
\begin{proof}
	Como $\Omega_x$ é enumerável, então presuma que $\Omega_X = \{\alpha_1, \ldots, \alpha_n\}$
	e que $\Parts{\Omega} = \{E_1, \ldots, E_n\}$ seja uma partição de $\Omega$ tal que
	\[
		\forall i \in [n],\ (P(E_i) = p_X(\alpha_i))
	\]
	Com isso,
	\[
		P\left(\bigcup_{i=1}^n E_i\right) =
		\sum_{i = 1}^n P(E_i) =
		\sum_{i = 1}^n p_X(\alpha_i)
	\]
	donde se conclui
	\[
		\sum_{i = 1}^n p_X(\alpha_i) = 1
	\]
\end{proof}
\end{prop:fmp suporte}
